\subsubsection{Répartition des paires entre processeur et accélérateur}\label{sec:routage-matrice}
Sur les fenêtres chromosomiques, WFA2 progresse rapidement dans les régions identiques. Le temps de NvPD sur GPU dépend surtout des dimensions de la grille. Pour choisir le solveur avant l’alignement, le programme estime la ressemblance de deux séquences à partir des mots de quinze bases qu’elles partagent.

Soit \(K_i\) l’ensemble des mots distincts de quinze bases A, C, G ou T lus sur le brin fourni de la séquence \(i\). Un symbole ambigu interrompt le mot en cours. Suivant le principe d’échantillonnage de Mash~\cite{mash}, le programme conserve les 256 plus petits hachages distincts de chaque ensemble. La fusion de deux listes donne l’échantillon \(S_{ij}\) des plus petits hachages de leur union. Si \(H_i\) désigne les hachages de \(K_i\), la proportion commune estime l’indice de Jaccard,
\[
  \widehat J_{ij}=\frac{|S_{ij}\cap H_i\cap H_j|}{|S_{ij}|},
  \qquad |S_{ij}|=\min(256,|H_i\cup H_j|).
\]
Avec deux listes non vides, une valeur d’au moins \(0{,}25\) envoie la paire vers WFA2 sur CPU. En dessous de ce seuil, la paire rejoint la file de NvPD sur GPU. Une liste vide déclenche le calcul par Edlib. Dans la variante avec reprise, chaque fil CPU qui a terminé les paires proches peut retirer une paire de l’autre extrémité de la file du GPU et la calculer avec Edlib.

La figure~\ref{fig:routage-matrice} compare sept façons de calculer la matrice sur un GH200 doté de quatre cœurs Grace réservés. Trois lots de huit fenêtres de 100\,000 bases proviennent du chromosome~21 de GRCh38 et de CHM13v2.0~\cite{grch38,t2t,ucscchain}. Le premier lot couvre un locus, le deuxième en couvre deux et le dernier huit. Chaque paire fournit deux entrées symétriques de la matrice. Les témoins exécutent toutes les paires avec le même solveur, Edlib ou WFA2 sur quatre cœurs, ou NvPD sur GPU. Une autre variante conserve le seuil de ressemblance et confie les paires distantes à Edlib sur CPU.

Le chronométrage commence après lecture du FASTA, avant l’allocation de la matrice et la construction des esquisses. Il s’arrête quand tous les scores sont en mémoire. À chaque bloc, chaque politique effectue un calcul d’échauffement, puis un calcul chronométré. Au fil des sept blocs, chaque politique occupe une fois chacune des sept positions dans l’ordre de passage.
\chart{comparison/routage_matrice.png}{Le panneau a) donne les durées moyennes pour 28 paires, avec l’écart-type et l’intervalle de Student à 95\,\% sur sept blocs. Le panneau b) compte les paires traitées par NvPD sur GPU. Le partage fixe lui attribue une paire sur deux.}{fig:routage-matrice}{python/studies/hybrid_matrix.py}

Au locus unique, les 28 paires restent sur CPU avec WFA2. Sur les deux loci, l’esquisse sépare 12 paires proches (31 à 1\,336 éditions) de 16 paires distantes (52\,184 à 52\,337 éditions). Avec reprise, NvPD traite huit paires distantes et les fils CPU reprennent les huit autres avec Edlib. La matrice prend \(0{,}990\pm0{,}005\)\,s, contre \(1{,}342\pm0{,}005\)\,s avec Edlib seul sur quatre cœurs et \(2{,}474\pm0{,}016\)\,s avec NvPD seul sur GPU. L’estimation et la mise en file prennent \(0{,}673\pm0{,}021\)\,ms, incluses dans le temps total.

Sur huit loci, les 28 paires sont distantes. La reprise et le partage fixe en attribuent chacun 14 au GPU. Le rapport apparié du temps fixe au temps avec reprise vaut 0,997 (IC à 95\,\% de 0,990 à 1,004). Sur la cohorte virale, la section~\ref{sec:routage-composition} examine le cas où toutes les paires passent le seuil de ressemblance.

\synthese{Sur le lot de deux loci, la reprise réduit le temps de 19,9\,\% par rapport au partage fixe (IC apparié à 95\,\% de 19,1 à 20,7\,\%). Sur huit loci, les deux répartitions aboutissent à la même charge de 14 paires par dispositif.}
